\chapter{Modelagem Matemática}
\label{cap:02}


\section{Introdução}

\textbf{2.1 Introdução à Modelagem Matemática.}

Objetivos de Aprendizagem

Compreender o que são modelos matemáticos, Aprender a construir modelos para sistemas biológicos, Dominar técnicas de análise de modelos e Aplicar modelagem a problemas reais.

\textbf{Definição (Modelo Matemático).}

Representação abstrata de um sistema real usando linguagem matemática, que permite fazer previsões quantitativas sobre o comportamento do sistema.

\textbf{Por que Modelar Sistemas Biológicos?.}

Previsão de comportamento, Teste de hipóteses in silico, Otimização de experimentos, Identificação de componentes chave e Design racional de sistemas.

Dinâmica de populações, Vias metabólicas, Redes regulatórias, Farmacocinética, Biologia sintética e Medicina personalizada.

--- IMP: 

"All models are wrong, but some are useful" --- George Box

\textbf{O Ciclo da Modelagem.}

%
\begin{figure}[htbp]
\centering
\begin{tikzpicture}[
auto,
    excalidraw/.style={
        decorate,
        decoration={random steps, segment length=3pt, amplitude=0.5pt},
        thick,
        font=\fontfamily{pzc}\selectfont\large
    },
    box/.style={
        rectangle, 
        draw=azulescuro, 
        excalidraw,
        fill=azulclaro!5, 
        text width=2.4cm, 
        text centered, 
        minimum height=1.2cm
    },
    arrow/.style={
        ->, 
        thick, 
        azulescuro, 
        excalidraw
    }
]

    \node[box] (bio) {Sistema Biológico};
    \node[box, right=2cm of bio] (modelo) {Modelo Matemático};
    \node[box, below=1.5cm of modelo] (analise) {Analise Computacional};
    \node[box, left=2cm of analise] (pred) {Previsões\\Hipóteses};

    \draw[arrow] (bio) -- node[above, font=\footnotesize] {Abstração} (modelo);
    \draw[arrow] (modelo) -- node[right, font=\footnotesize] {Simulação} (analise);
    \draw[arrow] (analise) -- node[below, font=\footnotesize] {Interpretação} (pred);
    \draw[arrow] (pred) -- node[left, font=\footnotesize] {Validação} (bio);
\end{tikzpicture}
\caption{Diagrama esquemático.}
\end{figure}
%

Ciclo Iterativo

O processo é iterativo: modelos são refinados com base em novos dados experimentais

\section{Tipos de Modelos}

\textbf{Classificação de Modelos.}

Por Natureza Matemática

resultado único para cada conjunto de condições e incorporam aleatoriedade e probabilidade.

Por Dependência Temporal

não mudam com o tempo (redes, equilíbrio) e evolução temporal (EDOs, EDPs).

Por Escala

nível molecular/celular, populações celulares e tecidos, órgãos, organismos.

\textbf{Modelos Contínuos vs Discretos.}

Modelos Contínuos

Variáveis contínuas, Equações diferenciais, Concentrações e Grande número de moléculas.

%
\begin{equation*}
\derivada{[X]}{t} = k_1 - k_2[X]
\end{equation*}
%

Modelos Discretos

Contagem de moléculas, Algoritmos estocásticos, Eventos discretos e Pequeno número de moléculas.

%
\begin{equation*}
X(t+1) = X(t) + \text{prod} - \text{deg}
\end{equation*}
%

Quando usar cada tipo?

Contínuo: 100 moléculas e Discreto: 100 moléculas.

\textbf{Modelos Empíricos vs Mecanísticos.}

Empíricos

Ajuste de curvas, Regressão estatística, Machine learning e Pouca interpretação biológica.

Simples e rápidos e Bom ajuste aos dados.

Pouco poder preditivo e Extrapolação arriscada.

Mecanísticos

Representam processos reais, Leis de conservação, Princípios físico-químicos e Interpretação biológica.

Poder preditivo e Extrapolação confiável.

Complexos e Muitos parâmetros.

\section{Construção de Modelos}

\textbf{Etapas da Construção de Modelos.}

\textbf{-}

Objetivo claro, Escopo bem definido e Questões a responder.

\textbf{-}

Variáveis de estado, Parâmetros e Entradas e saídas.

\textbf{-}

Leis de conservação, Cinéticas de reação e Condições iniciais e de contorno.

\textbf{-}

\textbf{Exemplo: Modelo de Decaimento Radioativo.}

Sistema

Um isótopo radioativo decai espontaneamente ao longo do tempo

 $N(t)$ = número de núcleos no tempo $t$

 Taxa de decaimento proporcional ao número de núcleos

--- DEST: 

%
\begin{equation*}
\derivada{N}{t} = -\lambda N
\end{equation*}
%

%
\begin{equation*}
N(t) = N_0 e^{-\lambda t}
\end{equation*}
%

onde $N_0 = N(0)$ é a condição inicial e $\lambda$ é a constante de decaimento.

\textbf{Implementação: Decaimento Radioativo.}

%
\begin{lstlisting}[language=Python, caption=Solução numérica do decaimento, basicstyle=\ttfamily\footnotesize]
import numpy as np
import matplotlib.pyplot as plt
from scipy.integrate import odeint

# Definir a EDO
def decay(N, t, lam):
    """Taxa de variacao de N"""
    return -lam * N

# Parametros
lambda_val = 0.5  # constante de decaimento
N0 = 100          # quantidade inicial
t = np.linspace(0, 10, 100)
\end{lstlisting}
%

\textbf{Implementação: Decaimento Radioativo (2/2).}

%
\begin{lstlisting}[language=Python, caption=Solução e Visualização]
# Resolver numericamente
N_num = odeint(decay, N0, t, args=(lambda_val,))

# Solucao analitica
N_ana = N0 * np.exp(-lambda_val * t)

# Plotar
plt.plot(t, N_num, 'b-', label='Numerica')
plt.plot(t, N_ana, 'r--', label='Analitica')
plt.xlabel('Tempo')
plt.ylabel('N(t)')
plt.legend()
\end{lstlisting}
%

\section{Equações Diferenciais em Biologia}

\textbf{Equações Diferenciais Ordinárias (EDOs).}

\textbf{Definição (EDO).}

Equação que relaciona uma função com suas derivadas em relação a uma variável independente (geralmente o tempo).

Forma Geral

%
\begin{equation*}
\derivada{\vect{x}}{t} = \vect{f}(\vect{x}, t, \vect{p})
\end{equation*}
%

onde:

$\vect{x} = $ vetor de variáveis de estado, $t = $ tempo, $\vect{p} = $ vetor de parâmetros e $\vect{f} = $ função que define a dinâmica.

\textbf{Classificação de EDOs.}

Por Ordem

$t = f(x,t)$, $dt^2 = f(x, t, t)$ e derivadas até $n$-ésima ordem.

Por Linearidade

$t = ax + b$ e $t = ax^2 + bx$.

Por Autonomia

$t = f(x)$ (não depende explicitamente de $t$) e $t = f(x,t)$.

\textbf{Exemplo: Crescimento Exponencial.}

Modelo Simples de Crescimento

Uma populacao cresce proporcionalmente ao seu tamanho:

%
\begin{equation*}
\derivada{N}{t} = rN
\end{equation*}
%

onde $r$ é a taxa de crescimento per capita.

%
\begin{equation*}
N(t) = N_0 e^{rt}
\end{equation*}
%

 crescimento\\

 decaimento\\

 estacionário

Tempo de Duplicação

$t_2 = r$

\textbf{Exemplo: Crescimento Logístico.}

Modelo de Verhulst

Populacao com capacidade de suporte limitada $K$:

%
\begin{equation*}
\derivada{N}{t} = rN\left(1 - \frac{N}{K}\right)
\end{equation*}
%

%
\begin{equation*}
N(t) = \frac{K N_0 e^{rt}}{K + N_0(e^{rt} - 1)}
\end{equation*}
%

$N K$ quando $t $, Forma sigmoidal e Ponto de inflexão em $N = K/2$.

%
\begin{figure}[htbp]
\centering
\begin{tikzpicture}
\draw[->] (0,0) -- (5,0) node[right] {$t$};
    \draw[->] (0,0) -- (0,4) node[above] {$N$};
    \draw[thick, azulescuro, domain=0:4.5, smooth, samples=100]
        plot (\x, {3/(1+2*exp(-1.5*\x))});
    \draw[dashed, verdeacido] (0,3) -- (5,3) node[right, font=\tiny] {$K$};
    \node[below] at (2.5,0) {Tempo};
\end{tikzpicture}
\caption{Diagrama esquemático.}
\end{figure}
%

\textbf{Implementação: Crescimento Logístico.}

%
\begin{lstlisting}[language=Python, caption=Modelo logístico, basicstyle=\ttfamily\footnotesize]
import numpy as np
import matplotlib.pyplot as plt
from scipy.integrate import odeint

def logistic(N, t, r, K):
    """Modelo de crescimento logistico"""
    return r * N * (1 - N/K)

# Parametros
r = 0.5    # taxa de crescimento
K = 100    # capacidade de suporte
N0 = 10    # populacao inicial

# Tempo
t = np.linspace(0, 20, 200)

# Resolver
N = odeint(logistic, N0, t, args=(r, K))
\end{lstlisting}
%

\textbf{Implementação: Crescimento Logístico (2/2).}

%
\begin{lstlisting}[language=Python, caption=Visualização]
# Plotar
plt.figure(figsize=(10, 6))
plt.plot(t, N, 'b-', linewidth=2, label='N(t)')
plt.axhline(y=K, color='r', linestyle='--',
            label=f'K = {K}')
plt.xlabel('Tempo')
plt.ylabel('Populacao N(t)')
plt.legend()
plt.grid(True)
\end{lstlisting}
%

\section{Sistemas de EDOs}

\textbf{Sistemas Acoplados.}

Sistema Geral

Quando múltiplas variáveis interagem:

%
\begin{align*}
\derivada{x_1}{t} &= f_1(x_1, x_2, \ldots, x_n, t)\\
\derivada{x_2}{t} &= f_2(x_1, x_2, \ldots, x_n, t)\\
&\vdots\\
\derivada{x_n}{t} &= f_n(x_1, x_2, \ldots, x_n, t)
\end{align*}
%

--- DEST: 

%
\begin{equation*}
\derivada{\vect{x}}{t} = \vect{f}(\vect{x}, t)
\end{equation*}
%

\textbf{Exemplo: Modelo Predador-Presa (Lotka-Volterra).}

Sistema

%
\begin{align*}
\derivada{V}{t} &= \alpha V - \beta V P \quad \text{(Presas)}\\
\derivada{P}{t} &= \delta \beta V P - \gamma P \quad \text{(Predadores)}
\end{align*}
%

onde:

$V$ = populacao de presas, $P$ = populacao de predadores, $$ = taxa de crescimento das presas, $$ = taxa de predação, $$ = taxa de morte dos predadores e $$ = eficiência de conversão.

\textbf{Interpretação do Modelo Predador-Presa.}

%
\begin{equation*}
\derivada{V}{t} = \underbrace{\alpha V}_{\text{crescimento}} - \underbrace{\beta V P}_{\text{predação}}
\end{equation*}
%

Crescem exponencialmente sem predadores, Perdem individuos por predação e Taxa de predação $\beta$ encontros ($VP$).

%
\begin{equation*}
\derivada{P}{t} = \underbrace{\delta \beta V P}_{\text{crescimento}} - \underbrace{\gamma P}_{\text{morte}}
\end{equation*}
%

Crescem ao consumir presas, Morrem naturalmente e Dependem das presas para sobreviver.

--- IMP: 

Este sistema exibe oscilações periódicas características!

\textbf{Implementação: Lotka-Volterra.}

%
\begin{lstlisting}[language=Python, caption=Sistema predador-presa, basicstyle=\ttfamily\footnotesize]
import numpy as np
import matplotlib.pyplot as plt
from scipy.integrate import odeint

def lotka_volterra(y, t, alpha, beta, delta, gamma):
    """Sistema Lotka-Volterra"""
    V, P = y
    dVdt = alpha*V - beta*V*P
    dPdt = delta*beta*V*P - gamma*P
    return [dVdt, dPdt]

# Parametros
alpha, beta, delta, gamma = 1.0, 0.1, 0.75, 1.5
y0, t = [10, 5], np.linspace(0, 50, 1000)
\end{lstlisting}
%

\textbf{Implementação: Lotka-Volterra (2/2).}

%
\begin{lstlisting}[language=Python, caption=Resolução e Plotting]
# Resolver
sol = odeint(lotka_volterra, y0, t,
             args=(alpha, beta, delta, gamma))

# Plotar series temporais
plt.plot(t, sol[:, 0], 'b-', label='Presas')
plt.plot(t, sol[:, 1], 'r-', label='Predadores')
plt.xlabel('Tempo')
plt.ylabel('Populacao')
plt.legend()
\end{lstlisting}
%

\textbf{Retrato de Fase.}

%
\begin{lstlisting}[language=Python, caption=Plano de fase, basicstyle=\ttfamily\tiny]
# Retrato de fase
plt.figure()
plt.plot(sol[:, 0], sol[:, 1],
         'g-', linewidth=2)
plt.plot(sol[0, 0], sol[0, 1],
         'go', markersize=10,
         label='Inicio')
plt.xlabel('Presas (V)')
plt.ylabel('Predadores (P)')
plt.title('Retrato de Fase')
plt.legend()
plt.grid(True)

# Campo vetorial
V = np.linspace(0, 20, 20)
P = np.linspace(0, 15, 20)
V_grid, P_grid = np.meshgrid(V, P)
dV, dP = lotka_volterra(
    [V_grid, P_grid], 0,
    alpha, beta, delta, gamma)
plt.quiver(V_grid, P_grid,
           dV, dP, alpha=0.5)
\end{lstlisting}
%

%
\begin{figure}[htbp]
\centering
\begin{tikzpicture}
\draw[->] (0,0) -- (4,0) node[right] {$V$};
    \draw[->] (0,0) -- (0,4) node[above] {$P$};

    % Trajetoria em espiral
    \draw[thick, verdeacido, domain=0:720, samples=200, smooth]
        plot ({2 + 1.5*cos(\x)/(1+\x/360)},
              {2 + 1.5*sin(\x)/(1+\x/360)});

    % Ponto de equilibrio
    \fill[red] (2,2) circle (2pt) node[right, font=\tiny] {Equilíbrio};

    \node[below] at (2,-0.2) {Presas};
    \node[left, rotate=90] at (-0.3,2) {Predadores};
\end{tikzpicture}
\caption{Diagrama esquemático.}
\end{figure}
%

Trajetórias fechadas, Oscilações periódicas e Centro estável.

\section{Analise de Equilíbrio}

\textbf{Pontos de Equilíbrio (Steady States).}

\textbf{Definição (Ponto de Equilíbrio).}

Estado no qual o sistema não muda com o tempo:

%
\begin{equation*}
\derivada{\vect{x}}{t} = \vect{0} \quad \Rightarrow \quad \vect{f}(\vect{x}^*, t) = \vect{0}
\end{equation*}
%

Exemplo: Crescimento Logístico

%
\begin{equation*}
\derivada{N}{t} = rN\left(1 - \frac{N}{K}\right) = 0
\end{equation*}
%

Pontos de equilíbrio:

$N^*_1 = 0$ (extinção) e $N^*_2 = K$ (capacidade de suporte).

\textbf{Estabilidade de Pontos de Equilíbrio.}

Classificação

sistema retorna ao equilíbrio após perturbação, sistema se afasta após perturbação e comportamento intermediário.

%
\begin{figure}[htbp]
\centering
\begin{tikzpicture}
% Estavel
    \begin{scope}
        \draw[->] (-1,0) -- (3,0) node[right] {$x$};
        \draw[->] (0,-1) -- (0,2) node[above] {$\dot{x}$};
        \draw[thick, azulescuro] (-0.5,1.5) parabola bend (1,0) (2.5,1.5);
        \fill[verdeacido] (1,0) circle (2pt) node[below, font=\tiny] {Estável};
        \draw[->, thick, red] (0.5,0.3) -- (1,0.1);
        \draw[->, thick, red] (1.5,0.3) -- (1,0.1);
    \end{scope}

    % Instavel
    \begin{scope}[xshift=5cm]
        \draw[->] (-1,0) -- (3,0) node[right] {$x$};
        \draw[->] (0,-1) -- (0,2) node[above] {$\dot{x}$};
        \draw[thick, azulescuro] (-0.5,-1.5) parabola[bend at end] (1,0);
        \draw[thick, azulescuro] (1,0) parabola (2.5,-1.5);
        \fill[red] (1,0) circle (2pt) node[below, font=\tiny] {Instável};
        \draw[->, thick, verdeacido] (0.5,-0.3) -- (0,-0.5);
        \draw[->, thick, verdeacido] (1.5,-0.3) -- (2,-0.5);
    \end{scope}
\end{tikzpicture}
\caption{Diagrama esquemático.}
\end{figure}
%

\textbf{Analise de Estabilidade Linear.}

Linearização em Torno do Equilíbrio

Para $ = ^* + $ (pequena perturbação):

%
\begin{equation*}
\derivada{\delta\vect{x}}{t} = \mathbf{J}(\vect{x}^*) \delta\vect{x}
\end{equation*}
%

onde $\mathbf{J}(\vect{x}^*)$ é a matriz Jacobiana:

%
\begin{equation*}
J_{ij} = \parcial{f_i}{x_j}\Bigg|_{\vect{x}^*}
\end{equation*}
%

Critério de Estabilidade

O equilíbrio $^*$ é estável se todos os autovalores de $(^*)$ têm parte real negativa.

\textbf{Exemplo: Analise do Crescimento Logístico.}

Modelo

%
\begin{equation*}
\derivada{N}{t} = f(N) = rN\left(1 - \frac{N}{K}\right)
\end{equation*}
%

%
\begin{equation*}
J = \derivada{f}{N} = r\left(1 - \frac{2N}{K}\right)
\end{equation*}
%

$N^* = 0$: $J(0) = r > 0$ $\Rightarrow$ instável e $N^* = K$: $J(K) = -r < 0$ $\Rightarrow$ estável.

--- IMP: 

A populacao converge para a capacidade de suporte $K$

\textbf{Analise de Estabilidade Computacional.}

%
\begin{lstlisting}[language=Python, caption=Calculo do Jacobiano e autovalores, basicstyle=\ttfamily\footnotesize]
import numpy as np
from scipy.linalg import eig

def jacobian_lotka_volterra(y, alpha, beta, delta, gamma):
    V, P = y
    return np.array([
        [alpha - beta*P, -beta*V],
        [delta*beta*P, delta*beta*V - gamma]
    ])

# Ponto de equilibrio e Jacobiano
V_eq, P_eq = gamma / (delta * beta), alpha / beta
J = jacobian_lotka_volterra([V_eq, P_eq], alpha, beta, delta, gamma)

# Autovalores
eigenvalues, _ = eig(J)
print(f"Autovalores: {eigenvalues}")
\end{lstlisting}
%

\textbf{Analise de Estabilidade Computacional.}

%
\begin{lstlisting}[language=Python, caption=Analise de autovalores]
# Ponto de equilibrio
V_eq = gamma / (delta * beta)
P_eq = alpha / beta

# Calcular Jacobiano no equilibrio
J = jacobian_lotka_volterra([V_eq, P_eq],
                             alpha, beta, delta, gamma)

# Autovalores
eigenvalues, eigenvectors = eig(J)
print(f"Autovalores: {eigenvalues}")
print(f"Partes reais: {eigenvalues.real}")
\end{lstlisting}
%

\section{Métodos Numéricos}

\textbf{Por que Métodos Numéricos?.}

Limitações das Soluções Analíticas

Maioria dos sistemas biológicos são não-lineares, Poucos modelos têm solução analítica fechada e Sistemas complexos requerem aproximação numérica.

Métodos Numéricos Comuns

simples, primeira ordem, preciso, quarta ordem, passo variável (RK45, Dormand-Prince) e sistemas rígidos (BDF).

\textbf{Método de Euler.}

Ideia Básica

Aproximar derivada por diferença finita:

%
\begin{equation*}
\derivada{x}{t} \approx \frac{x(t+h) - x(t)}{h}
\end{equation*}
%

--- DEST: 

%
\begin{equation*}
x(t+h) = x(t) + h \cdot f(x(t), t)
\end{equation*}
%

Partir de $x_0$ em $t_0$, Calcular $x_1 = x_0 + h f(x_0, t_0)$ e Repetir: $x_n+1 = x_n + h f(x_n, t_n)$.

Erro

Erro local: $O(h^2)$, Erro global: $O(h)$

\textbf{Implementação do Método de Euler.}

%
\begin{lstlisting}[language=Python, caption=Euler forward, basicstyle=\ttfamily\footnotesize]
def euler(f, x0, t):
    n = len(t)
    x = np.zeros((n, len(x0)))
    x[0] = x0
    for i in range(n-1):
        h = t[i+1] - t[i]
        x[i+1] = x[i] + h * np.array(f(x[i], t[i]))
    return x

# Exemplo
def exp_growth(x, t, r=0.5): return r * x
t = np.linspace(0, 10, 100)
x_euler = euler(lambda x,t: exp_growth(x,t), [1.0], t)
\end{lstlisting}
%

\textbf{Método de Runge-Kutta de 4ª Ordem (RK4).}

Ideia

Usar múltiplas avaliações da derivada para melhor aproximação

--- DEST: 

%
\begin{align*}
k_1 &= f(x_n, t_n)\\
k_2 &= f(x_n + \frac{h}{2}k_1, t_n + \frac{h}{2})\\
k_3 &= f(x_n + \frac{h}{2}k_2, t_n + \frac{h}{2})\\
k_4 &= f(x_n + hk_3, t_n + h)\\
x_{n+1} &= x_n + \frac{h}{6}(k_1 + 2k_2 + 2k_3 + k_4)
\end{align*}
%

Precisão

Erro local: $O(h^5)$, Erro global: $O(h^4)$

\textbf{Implementação do RK4.}

%
\begin{lstlisting}[language=Python, caption=Runge-Kutta 4, basicstyle=\ttfamily\footnotesize]
def rk4(f, x0, t):
    n, x = len(t), np.zeros((len(t), len(x0)))
    x[0] = x0
    for i in range(n-1):
        h = t[i+1] - t[i]
        k1 = np.array(f(x[i], t[i]))
        k2 = np.array(f(x[i] + h*k1/2, t[i] + h/2))
        k3 = np.array(f(x[i] + h*k2/2, t[i] + h/2))
        k4 = np.array(f(x[i] + h*k3, t[i] + h))
        x[i+1] = x[i] + (h/6) * (k1 + 2*k2 + 2*k3 + k4)
    return x

# Usar
x_rk4 = rk4(lambda x,t: exp_growth(x,t), [1.0], t)
\end{lstlisting}
%

\section{Analise de Sensibilidade}

\textbf{Analise de Sensibilidade.}

\textbf{Definição (Sensibilidade).}

Medida de como a saída de um modelo varia em resposta a mudanças nos parâmetros de entrada.

Por que é importante?

Identificar parâmetros críticos, Priorizar experimentos para estimar parâmetros, Avaliar robustez do modelo e Simplificar modelos complexos.

Coeficiente de Sensibilidade

%
\begin{equation*}
S_{x,p} = \parcial{x}{p} \cdot \frac{p}{x} = \frac{\parcial \ln x}{\parcial \ln p}
\end{equation*}
%

Sensibilidade relativa normalizada

\textbf{Tipos de Analise de Sensibilidade.}

Local

Variação em torno de um ponto nominal, Derivadas parciais e Rápida e analítica (quando possível).

Global

Exploração de todo o espaço de parâmetros, Métodos de amostragem (Monte Carlo, Sobol) e Computacionalmente intensiva.

--- Trade-off ---

Local: rápida mas limitada \\

Global: abrangente mas custosa

\textbf{Analise de Sensibilidade Local.}

%
\begin{lstlisting}[language=Python, caption=Sensibilidade por diferenças finitas, basicstyle=\ttfamily\footnotesize]
def sensitivity_local(model, params, param_names,
                       perturbation=0.01):
    """
    Analise de sensibilidade local
    """
    baseline = model(params)
    sensitivities = {}
    for i, name in enumerate(param_names):
        params_pert = params.copy()
        params_pert[i] *= (1 + perturbation)
        output_pert = model(params_pert)
        sensitivities[name] = (output_pert - baseline) / baseline / perturbation
    return sensitivities

# Exemplo
sens = sensitivity_local(lambda p: p[1], [0.5, 100, 10], ['r', 'K', 'N0'])
\end{lstlisting}
%

\section{Ajuste de Modelos a Dados}

\textbf{Estimação de Parâmetros.}

Problema

Dados: $(t_1, y_1), (t_2, y_2), , (t_n, y_n)$\\

Modelo: $(t; )$\\

Objetivo: Encontrar $^*$ que melhor ajusta os dados

Função Objetivo

Mínimos quadrados:

%
\begin{equation*}
\text{SSE}(\vect{p}) = \sum_{i=1}^{n} [y_i - \hat{y}(t_i; \vect{p})]^2
\end{equation*}
%

Objetivo: $\min_ ()$

Métodos de Otimização

Levenberg-Marquardt, Evolução Diferencial e Algoritmos Genéticos.

\textbf{Ajuste de Curvas com scipy.}

%
\begin{lstlisting}[language=Python, caption=Estimação de parâmetros, basicstyle=\ttfamily\footnotesize]
from scipy.optimize import curve_fit

# Dados experimentais
t_data = np.linspace(0, 10, 20)
N_data = 100/(1+9*np.exp(-0.5*t_data)) + np.random.normal(0, 2, 20)

# Modelo e Ajuste
def logistic_model(t, r, K): return K / (1 + (K/10 - 1)*np.exp(-r*t))
p_opt, _ = curve_fit(logistic_model, t_data, N_data, p0=[1.0, 50])

# Plotar
plt.plot(t_data, N_data, 'o', label='Dados')
plt.plot(t_data, logistic_model(t_data, *p_opt), '-', label='Ajuste')
\end{lstlisting}
%

\textbf{Validação de Modelos.}

Métricas de Qualidade do Ajuste

Soma dos erros quadrados, Coeficiente de determinação (0 a 1), Raiz do erro quadrático médio e Critérios de informação (penalizam complexidade).

%
\begin{equation*}
R^2 = 1 - \frac{\sum (y_i - \hat{y}_i)^2}{\sum (y_i - \bar{y})^2}
\end{equation*}
%

Atenção!

$R^2$ alto não garante bom modelo, Validar com dados independentes e Verificar plausibilidade biológica.

\section{Estudo de Caso}

\textbf{Estudo de Caso: Dinâmica de Infecção Viral.}

Sistema

Modelo SIR (Susceptible-Infected-Recovered):

%
\begin{align*}
\derivada{S}{t} &= -\beta \frac{SI}{N}\\
\derivada{I}{t} &= \beta \frac{SI}{N} - \gamma I\\
\derivada{R}{t} &= \gamma I
\end{align*}
%

Número Básico de Reprodução ($R_0$)

%
\begin{equation*}
R_0 = \frac{\beta}{\gamma}
\end{equation*}
%

\textbf{Estudo de Caso: Dinâmica de Infecção Viral (cont.).}

Parâmetros do Modelo SIR

onde:

$S$ = susceptíveis, $I$ = infectados, $R$ = recuperados, $$ = taxa de transmissão, $$ = taxa de recuperação e $N = S + I + R$ = populacao total.

\textbf{Numero Básico de Reprodução.}

--- DEF: $R_0$ (R-naught) ---

Numero médio de infecções secundárias causadas por um indivíduo infectado em uma populacao totalmente susceptível.

--- DEST: 

%
\begin{equation*}
R_0 = \frac{\beta}{\gamma}
\end{equation*}
%

Interpretação

$R_0 < 1$: epidemia se extingue, $R_0 > 1$: epidemia se propaga e $R_0 = 1$: limiar epidêmico.

$R_0$ estimado entre 2-6, dependendo da variante e medidas de controle

\textbf{Implementação do Modelo SIR.}

%
\begin{lstlisting}[language=Python, caption=Modelo SIR completo, basicstyle=\ttfamily\footnotesize]
def sir_model(y, t, beta, gamma):
    S, I, R = y
    N = S + I + R
    return [-beta * S * I / N, beta * S * I / N - gamma * I, gamma * I]

# Parametros e condicoes iniciais
beta, gamma = 0.5, 0.1
R0 = beta / gamma
N, I0 = 1000, 10
y0 = [N - I0, I0, 0]
\end{lstlisting}
%

\textbf{Analise e Visualização.}

%
\begin{lstlisting}[language=Python, caption=Analise do modelo SIR, basicstyle=\ttfamily\footnotesize]
S, I, R = sol.T
plt.figure(figsize=(10, 6))
plt.plot(t, S, 'b-', label='Suscept.'), plt.plot(t, I, 'r-', label='Infect.')
plt.plot(t, R, 'g-', label='Recuper.')
plt.xlabel('Tempo'), plt.ylabel('Individuos'), plt.legend(), plt.grid(True)

# Estatisticas
peak_idx = np.argmax(I)
print(f"Pico: {I[peak_idx]:.0f} em {t[peak_idx]:.1f} dias")
print(f"Ataque final: {R[-1]/N*100:.1f}%")
\end{lstlisting}
%

\section{Modelos Estocásticos: Reptation do DNA}

\textbf{Por que Modelos Estocásticos?.}

Limitações dos Modelos Determinísticos

Assumem grande número de moléculas e variáveis contínuas, Ignoram flutuações térmicas e variabilidade molecular e Inadequados quando $N 100$ moléculas ou eventos raros.

Quando usar modelos estocásticos?

Movimentos de moléculas individuais, Sistemas com ruído térmico relevante (escala nanométrica), Polímeros em confinamento (gel, cromatina, poros) e Expressão gênica com poucos fatores de transcrição.

--- IMP: 

Flutuações não são "ruído a eliminar" — frequentemente carregam informação biológica.

\textbf{Movimento de Polímeros: O Modelo de Reptation.}

\textbf{Definição (Reptation).}

Mecanismo de movimento de um polímero longo (como o DNA) confinado em uma rede ou gel, proposto por P.G.\ de Gennes (1971). O polímero se move como uma cobra dentro de um tubo virtual formado pelos obstáculos ao seu redor.

O DNA não pode atravessar os obstáculos da rede, Cria um ``tubo'' ao longo de seu contorno, Move-se pela criação/destruição de alças nas extremidades e Movimento resultante: difusão ao longo do próprio eixo.

%
\begin{figure}[htbp]
\centering
\begin{tikzpicture}
% Obstáculos
    \foreach \x/\y in {0.5/0.5, 1.5/1.5, 2.5/0.8, 3.5/1.8, 0.8/2.5, 2.0/2.8, 3.2/0.3, 1.2/3.2, 3.8/2.5} {
        \fill[gray!50] (\x, \y) circle (3pt);
    }
    % Tubo / DNA
    \draw[thick, azulescuro, rounded corners=4pt]
        (0,1.0) .. controls (0.5,0.9) and (1.0,1.3) ..
        (1.5,1.1) .. controls (2.0,0.9) and (2.5,1.5) ..
        (3.0,1.4) .. controls (3.5,1.3) and (3.8,1.8) .. (4.0,1.6);
    % Setas de movimento
    \draw[->, thick, verdeacido] (3.8,1.7) -- (4.3,1.55);
    \draw[->, thick, red]        (0.2,1.05) -- (-0.3,1.1);
    \node[font=\tiny, azulescuro] at (2.0, 0.3) {DNA};
    \node[font=\tiny, gray]       at (3.5, 3.0) {gel};
\end{tikzpicture}
\caption{Diagrama esquemático.}
\end{figure}
%

\textbf{Formulação Estocástica do Reptation (1/2).}

Modelo de Caminhada Aleatória no Tubo

O DNA é representado por $N$ segmentos de comprimento $b$. A cada passo de tempo $ t$, cada extremidade:

Com prob.\ $2$: — adiciona um segmento à frente e Com prob.\ $2$: — remove um segmento da frente.

--- DEST: 

%
\begin{equation*}
D_c = \frac{b^2}{2\,\tau}
\end{equation*}
%

\textbf{Formulação Estocástica do Reptation (2/2).}

%
\begin{equation*}
D \propto \frac{D_c}{N^2} \sim \frac{1}{N^2}
\end{equation*}
%

Predição chave

O tempo de renovação do tubo escala como $\tau_ rep N^3$ — verificado experimentalmente por eletroforese em gel.

--- IMP: 

Quanto mais longo o DNA, mais lento o movimento — base da separação por tamanho em gel.

\textbf{Reptation em Campo Elétrico (Eletroforese).}

Biased Reptation (Duke, 1989)

Na presença de campo elétrico $E$, a probabilidade de avanço/recuo nas extremidades se torna assimétrica:

%
\begin{align*}
p_+ &= \frac{1}{2}\left(1 + \frac{qEb}{k_BT}\right) \quad \text{(avança)}\\
p_- &= \frac{1}{2}\left(1 - \frac{qEb}{k_BT}\right) \quad \text{(recua)}
\end{align*}
%

%
\begin{equation*}
\mu = \frac{v}{E} \propto \frac{1}{N}
\end{equation*}
%

Fragmentos menores: maior mobilidade, Base da separação por gel eletroforese e Válido para DNA (regime de reptation).

Limite do modelo

Para campos altos ($qEb k_BT$): DNA se alinha com o campo e a separação por tamanho se perde.

\textbf{Implementação: Simulação de Reptation.}

%
\begin{lstlisting}[language=Python, caption=Simulação estocástica de reptation, basicstyle=\ttfamily\footnotesize]
import numpy as np
import matplotlib.pyplot as plt

def simular_reptation(N, n_steps, E=0.0, kBT=1.0, q=1.0, b=1.0):
    """
    Simula reptation de DNA em campo eletrico E.
    N       : numero de segmentos
    n_steps : passos de Monte Carlo
    E       : campo eletrico (0 = difusao pura)
    """
    # Posicao do centro de massa ao longo do tubo
    pos_cm = np.zeros(n_steps)
    pos = 0.0  # deslocamento acumulado

    # Probabilidade de avanco
    p_mais = 0.5 * (1.0 + q * E * b / kBT)
    p_mais = np.clip(p_mais, 0.0, 1.0)

    for t in range(1, n_steps):
        # Sortear qual extremidade se move (frente ou tras)
        r = np.random.rand()
        if r < p_mais:
            pos += b / N   # deslocamento do CM: b/N por passo
        else:
            pos -= b / N
        pos_cm[t] = pos

    return pos_cm
\end{lstlisting}
%

\textbf{Implementação: Análise e Visualização.}

%
\begin{lstlisting}[language=Python, caption=Comparacao entre tamanhos de DNA, basicstyle=\ttfamily\footnotesize]
n_steps = 50000
tamanhos = [50, 100, 200]   # numero de segmentos
E = 0.1                      # campo eletrico fraco

fig, axes = plt.subplots(1, 2, figsize=(12, 4))

for N in tamanhos:
    traj = simular_reptation(N, n_steps, E=E)
    t = np.arange(n_steps)

    # Trajetoria
    axes[0].plot(t, traj, label=f'N={N}', alpha=0.7)

    # MSD (mean squared displacement)
    msd = np.cumsum((traj - traj[0])**2) / np.arange(1, n_steps+1)
    axes[1].loglog(t[1:], msd[1:], label=f'N={N}')

axes[0].set(xlabel='Passos', ylabel='Posicao CM', title='Trajetorias')
axes[1].set(xlabel='Passos', ylabel='MSD', title='MSD (escala log-log)')
for ax in axes:
    ax.legend(); ax.grid(True)
plt.tight_layout()
\end{lstlisting}
%

\textbf{Predições e Verificação Experimental.}

Escalonamentos do Modelo de Reptation

Aplicações Biotecnológicas

separação de fragmentos de DNA por tamanho, entender mobilidade de fragmentos, separação de cromossomos inteiros e translocação de DNA — extensão do modelo de reptation.

\section{Boas Práticas em Modelagem}

\textbf{Princípios de Modelagem.}

1. Simplicidade vs Realismo

Comece simples, adicione complexidade gradualmente, "Make things as simple as possible, but not simpler" --- Einstein e Balanceie entre descrição acurada e tratabilidade.

2. Validação

Compare com dados experimentais, Teste previsões em novos experimentos e Valide em múltiplas escalas/condições.

3. Documentação

Documente hipóteses claramente, Registre origem de parâmetros e Mantenha código reproduzível.

\textbf{Armadilhas Comuns.}

Erros Frequentes

modelo complexo demais para os dados, múltiplas soluções, sempre verifique dimensões, valores negativos, etc, além do regime validado e sempre quantifique incerteza.

--- IMP: 

"Remember that all models are wrong; the practical question is how wrong do they have to be to not be useful" --- George Box

\textbf{Ferramentas Computacionais.}

Python

NumPy, SciPy, Matplotlib, Seaborn, PyDSTool, Tellurium e COBRApy (metabolismo).

R

deSolve, FME (fitting) e ggplot2.

MATLAB

ode45, ode15s, SimBiology e Simbólico.

Especializadas

COPASI, CellDesigner, Virtual Cell e BioNetGen.

\section{Exercícios}

\textbf{Exercícios - Parte 1.}

Questão 1: Modelo de Crescimento

Considere um modelo de crescimento com colheita constante:

%
\begin{equation*}
\derivada{N}{t} = rN\left(1 - \frac{N}{K}\right) - H
\end{equation*}
%

Encontre os pontos de equilíbrio e Analise a estabilidade para $H < rK/4$ e $H > rK/4$.

\textbf{Exercícios - Parte 2.}

Questão 2: Método Numérico

Implemente o método de Euler para resolver:

%
\begin{equation*}
\derivada{y}{t} = -y + \sin(t), \quad y(0) = 1
\end{equation*}
%

Compare com a solução analítica para diferentes passos.

\textbf{Exercícios - Parte 2.}

Questão 3: Sistema Acoplado

Considere o sistema de competição:

%
\begin{align*}
\derivada{N_1}{t} &= r_1 N_1 \left(1 - \frac{N_1 + \alpha N_2}{K_1}\right)\\
\derivada{N_2}{t} &= r_2 N_2 \left(1 - \frac{N_2 + \beta N_1}{K_2}\right)
\end{align*}
%

Implemente o modelo em Python, Encontre os pontos de equilíbrio, Simule para diferentes valores de $$ e $$ e O que acontece quando $ > 1$ vs $ < 1$?.

\textbf{Exercícios - Parte 3.}

Projeto: Modelo SIR com Vacinação

Estenda o modelo SIR para incluir vacinação:

%
\begin{align*}
\derivada{S}{t} &= -\beta \frac{SI}{N} - \nu S, \quad \derivada{I}{t} = \beta \frac{SI}{N} - \gamma I, \quad \derivada{R}{t} = \gamma I + \nu S
\end{align*}
%

Implemente o modelo e calcule o $R_0$ efetivo, Determine $\nu$ necessária para controle ($R_0^{\text{eff}} < 1$) e Simule diferentes estratégias de vacinação.

\section{Notas Bibliográficas}

\textbf{Referências Principais.}

99

 Strogatz S.H. (2018). Nonlinear Dynamics and Chaos, 2nd ed. CRC Press.

 Edelstein-Keshet L. (2005). Mathematical Models in Biology. SIAM.

 Britton N.F. (2003). Essential Mathematical Biology. Springer.

 Murray J.D. (2002). Mathematical Biology I: An Introduction, 3rd ed. Springer.

 Keener J., Sneyd J. (2009). Mathematical Physiology, 2nd ed. Springer.

\textbf{Leitura Complementar.}

Livros Avançados

Ellner S.P., Guckenheimer J. (2011). Dynamic Models in Biology, Alon U. (2019). An Introduction to Systems Biology, 2nd ed e Klipp E. et al. (2016). Systems Biology, 2nd ed.

Recursos Online

\textbf{-}

Mathematical Biology course.

\textbf{-}

Tutoriais Python

\textbf{-}

\textbf{-}

 Obrigado!

 Capítulo 2: Introdução à Modelagem Matemática

%

Static Network Models

Dúvidas? Perguntas?
